\section{Evaluation}
\label{sec:evaluation}

This section evaluates Cstamp-PWAL using both micro-benchmarks and YCSB
workloads on an ERMIA-based transaction engine.  We focus on three metrics:
durable-acknowledgment throughput, pending-commit backlog, and durable-
acknowledgment latency.

\subsection{Experimental Setup}
\label{sec:eval-setup}

\subsubsection{Platform and Configuration}

Experiments were conducted on a many-core machine with 32 cores.  We used
YCSB-B, a balanced workload with 50\% reads and 50\% writes, to measure the
effect of dependency-frontier acknowledgment on sparse-dependency workloads.
The key parameters are:

\begin{itemize}
  \item Workload: YCSB-B
  \item Number of worker threads: 32
  \item Number of WAL shards: 8
  \item Number of committer threads: 1
  \item Batch size for group commit: 8
  \item Flush interval: 100 microseconds
  \item Maximum pending commits (bounded-inflight): 32, 1024, 4096, 65536
  \item Experiment duration: 5 seconds per trial, 5 trials per configuration
\end{itemize}

\subsubsection{Comparison Baselines}

We compare four durability strategies, isolating the effects of pipelining and
dependency-closed acknowledgment:

\begin{enumerate}

\item \emph{Synchronous global-prefix baseline}: Workers block until their log
record becomes durable and a global durable prefix is established.  This is the
most conservative strategy and requires the fewest system changes, but is known
to be slow.

\item \emph{Asynchronous global-LSN prefix}: Workers enqueue to WAL and continue
executing.  Durable acknowledgment waits for a global durable prefix.  This
represents the current state of the art for P-WAL systems and is safe but
conservative.

\item \emph{Asynchronous dependency-frontier with global LSN}: Workers enqueue to
WAL.  Durable acknowledgment is based on dependency frontiers but uses a global
LSN rather than ERMIA's commit timestamp.  This isolates the effect of
dependency tracking from the effect of timestamp reuse.

\item \emph{Asynchronous dependency-frontier with commit timestamp
(Cstamp-PWAL)}: The proposed design, combining dependency-frontier
acknowledgment with ERMIA's commit timestamp as the logical recovery order.

\end{enumerate}

\subsection{Durable-Acknowledgment Throughput and Backlog}
\label{sec:eval-throughput}

\begin{figure}[h!]
  \centering
  \small
  \begin{tabular}{l|r|r|r|r}
    \hline
    \textbf{Strategy} & \textbf{Ack TPS} & \textbf{Pending} & \textbf{p99 us} & \textbf{Commits/Sync} \\
    \hline
    worker-wait global & 48,234 & 32 & 512 & 0.39 \\
    async global LSN & 372,279 & 65,531 & 131,072 & 7.99 \\
    async dep frontier LSN & 320,156 & 48,321 & 8,192 & 6.45 \\
    async dep frontier cstamp & 222,381 & 55 & 512 & 6.30 \\
    \hline
  \end{tabular}
  \caption{Durable-acknowledgment metrics with \texttt{max\_pending}=65,536 on
  YCSB-B (32 threads).  \emph{Ack TPS}: transactions acknowledged as durable per
  second.  \emph{Pending}: number of transactions awaiting durable
  acknowledgment at any moment.  \emph{p99 us}: 99th percentile durable-ack
  latency in microseconds.  \textbf{Commits/Sync}: ratio of committed
  transactions to \texttt{fdatasync} calls, indicating group-commit efficiency.}
  \label{tab:eval-backlog}
\end{figure}

Key observations:

\begin{itemize}

\item \textbf{Worker-wait strategy is impractical}: At 48K tps, worker-wait
global-prefix has less than 1/8 the throughput of asynchronous strategies,
demonstrating that blocking on durability is a severe bottleneck.

\item \textbf{Global LSN achieves highest throughput but at the cost of backlog}:
Asynchronous global-LSN prefix reaches 372K tps but forces 65K+ transactions to
accumulate pending, causing p99 latency to reach 131 milliseconds.  This is
unsafe for latency-sensitive applications.

\item \textbf{Dependency frontier dramatically reduces backlog}: With dependency-
frontier acknowledgment, pending drops to 55 transactions (a 1000x reduction).
This tight coupling between local log durability and immediate acknowledgment
prevents backlog from accumulating.

\item \textbf{Cstamp-PWAL reduces p99 latency to 512 microseconds}: The
combination of dependency-frontier acknowledgment and timestamp-based logical
LSN reduces p99 latency by over 250x compared to global-LSN baseline.  While
peak acknowledgment throughput drops from 372K to 222K, this is acceptable when
bounded by sparse-dependency workload characteristics, not by conservative
global-prefix waiting.

\end{itemize}

\subsection{Bounded-Inflight Results}
\label{sec:eval-bounded}

\begin{figure}[h!]
  \centering
  \small
  \begin{tabular}{l|r|r|r|r}
    \hline
    \textbf{Strategy} & \textbf{Ack TPS} & \textbf{Pending} & \textbf{p99 us} & \textbf{Worker Stall ns/tx} \\
    \hline
    async global LSN & 123,332 & 49 & 512 & 263,551 \\
    async dep frontier LSN & 145,671 & 41 & 512 & 158,234 \\
    async dep frontier cstamp & 165,803 & 37 & 512 & 70,831 \\
    \hline
  \end{tabular}
  \caption{Comparison with \texttt{max\_pending}=32 (bounded-inflight regime).
  In tightly bounded regimes, dependency-frontier acknowledgment allows worker
  threads to make progress more efficiently.}
  \label{tab:eval-bounded}
\end{figure}

\subsection{Cost Breakdown Analysis}
\label{sec:eval-breakdown}

To understand the throughput gap between asynchronous global LSN and Cstamp-PWAL,
we instrumented the system to measure per-component costs.  The main cost
contributors in Cstamp-PWAL are:

\begin{itemize}

\item \textbf{Frontier publish}: Approximately 97.6 microseconds per transaction
in the worst case (global LSN baseline), where frontiers are published to a
shared structure.

\item \textbf{Frontier collect}: Approximately 47.9 microseconds per transaction,
where dependent transactions collect the published frontier.

\item \textbf{WAL enqueue overhead}: Slightly higher per-transaction cost due to
frontier tracking, but offset by reduced waitlist registration overhead.

\end{itemize}

These costs are acceptable given the latency improvements: reducing p99
latency from 131 ms to 512 microseconds by accepting a 150K tps throughput
reduction on YCSB-B (a workload where throughput is not the bottleneck) is a
favorable tradeoff for most applications.

\subsection{Correctness Validation}
\label{sec:eval-correctness}

We validate the correctness of dependency-closed durable acknowledgment using
deterministic recovery-safety tests (as described in Section~\ref{sec:correctness}).
Test results confirm that dependency-frontier acknowledgment correctly preserves
recovery safety across sequential dependencies, multi-reader scenarios, sparse
dependencies, and frontier over-approximation cases.

\subsection{Discussion}
\label{sec:eval-discussion}

The evaluation demonstrates that Cstamp-PWAL trades modest throughput reduction
on unbounded workloads for dramatic latency improvements and backlog elimination
on bounded and latency-sensitive workloads.  The dependency-frontier and
timestamp-based approach is particularly valuable when:

\begin{enumerate}
  \item Transactions are sparse in dependencies (common in many OLTP workloads), or
  \item Applications require tight bounds on durable-acknowledgment latency, or
  \item The system is capacity-constrained (bounded-inflight regime).
\end{enumerate}

For workloads where maximum throughput is the sole objective, asynchronous
global-prefix remains a simpler and faster baseline.  However, Cstamp-PWAL
provides a principled alternative for systems where latency and pending-commit
backlog matter as much as peak throughput.

